% PAGES: 275-275
% Section 9.2 closed cleanly at the end of sec09_c1_8.tex. This file opens
% section 9.3 fresh.

\section{Minimum variance portfolios}

The minimum variance portfolio invested in $n$ assets and with a cash position is the
portfolio which has the smallest variance of its return of all the portfolios with a
given expected return $\mu_P$.

In other words, given $\mu_P$, we are looking for a weights vector $w = (w_i)_{i=1:n}$
such that $\mu_R = \mu_P$ and $\sigma_R^2$ is minimal. From (9.6) and (9.8), it follows
that we want to minimize $\sigma_R^2 = w^t\Sigma_Rw$ subject to $\mu_R = r_f +
\bar\mu^tw = \mu_P$, i.e.,
\begin{equation}
\min_{r_f+\bar\mu^tw=\mu_P} w^t\Sigma_Rw.
\end{equation}

We use Lagrange multipliers to solve this problem; see section 10.2.2 for details. The
associated Lagrangian function $F : \R^{n+1} \to \R$ is given by
\[
F(w,\lambda) \;=\; w^t\Sigma_Rw \;+\; \lambda(\mu_P-r_f-\bar\mu^tw),
\]
where $\lambda \in \R$ is the Lagrange multiplier.\footnote{Lagrange multipliers can be
used to solve this constrained optimization problem since the condition (10.66), i.e.,
$\rank(Dg(w)) = 1$, is satisfied for the constraint function $g(w) = \mu_P - r_f -
\bar\mu^tw$. From (10.43), we find that
\[
Dg(w) \;=\; D(\mu_P-r_f-\bar\mu^tw) \;=\; -D(\bar\mu^tw) \;=\; -\bar\mu^t,
\]
and therefore $\rank(Dg(w)) = 1$.}

Since the covariance matrix $\Sigma_R$ is symmetric, we obtain from (10.47) that
\begin{equation}
D\left(w^t\Sigma_Rw\right) \;=\; 2\left(\Sigma_Rw\right)^t.
% Printed with a second stray period immediately after this one, PDF page 275
% (folio 259) -- almost certainly scan/print noise rather than deliberate
% typesetting, consistent with other isolated ink specks on this page;
% verified at 600 dpi. Not reproduced as a second punctuation mark; see also
% the ch08 Put-Call-parity exercise note for this book's other stray-ink-dot
% artifact.
\end{equation}

Also, from (10.43), it follows that
\begin{equation}
D(\bar\mu^tw) \;=\; \bar\mu^t.
\end{equation}

Then,
\begin{align}
D_wF(w,\lambda) &= D\left(w^t\Sigma_Rw\right) + \lambda D(\mu_P-r_f-\bar\mu^tw) \notag\\
&= D\left(w^t\Sigma_Rw\right) - \lambda D(\bar\mu^tw) \notag\\
&= 2\left(\Sigma_Rw\right)^t - \lambda\bar\mu^t \notag\\
&= \left(2\Sigma_Rw-\lambda\bar\mu\right)^t.
\end{align}

Also, note that
\begin{equation}
D_\lambda F(w,\lambda) \;=\; \pd{F}{\lambda}(w,\lambda) \;=\; \mu_P-r_f-\bar\mu^tw.
\end{equation}

From (10.67), (9.24), and (9.25), we conclude that the gradient of the Lagrangian
function is
\[
D_{(w,\lambda)}F(w,\lambda) \;=\;
\begin{pmatrix} (D_wF(w,\lambda))^t \\ (D_\lambda F(w,\lambda))^t \end{pmatrix}^t
\;=\;
\begin{pmatrix} 2\Sigma_Rw-\lambda\bar\mu \\ \mu_P-r_f-\bar\mu^tw \end{pmatrix}^t.
\]

To find the critical point of the Lagrangian, we solve $D_{(w,\lambda)}F(w_0,\lambda_0)
= 0$, which is equivalent to solving
\begin{align}
2\Sigma_Rw_0 &= \lambda_0\bar\mu; \\
\bar\mu^tw_0 &= \mu_P-r_f.
\end{align}
